\SourceLeaf{20}
\par
We also note that, on substituting for $x,y,z$ their values expressed in terms of $p,q$, the function $W$ must necessarily be identically $=0$. Therefore the following must likewise hold identically, that is, independently of the values of $\mathrm{d}p,\mathrm{d}q$.
\[
0=(\lambda T+\mu U+\nu V)\mathrm{d}p+(\lambda'T+\mu'U+\nu'V)\mathrm{d}q
\]
Equivalently, we have
\[
\begin{aligned}
\lambda T+\mu U+\nu V&=0\\
\lambda'T+\mu'U+\nu'V&=0
\end{aligned}
\]
It follows that the quantities $\mu\nu'-\nu\mu'$, $\nu\lambda'-\lambda\nu'$, $\lambda\mu'-\mu\lambda'$ are respectively proportional to $T,U,V$, or to the cosines of the angles $QX,QY,QZ$. This could already have been inferred from what was said above, with the signs still undetermined.

\Article{11}
\par
We turn from these general investigations to elliptic spheroidal bodies, for whose sake they were undertaken. If the origin is taken at the center of the body and the semiaxes are denoted by $A,B,C$, the equation of the surface will be
\[
\frac{xx}{AA}+\frac{yy}{BB}+\frac{zz}{CC}=1
\]
\par
We therefore set $W=\frac{xx}{AA}+\frac{yy}{BB}+\frac{zz}{CC}-1$. It is clear that at every point inside the body $W$ takes a negative value, and at every point outside it a positive value. We further have $T=\frac{2x}{AA}$, $U=\frac{2y}{BB}$, $V=\frac{2z}{CC}$. Thus, setting
\[
\BareRoot{(\frac{xx}{A^4}+\frac{yy}{B^4}+\frac{zz}{C^4})}=\psi
\]
we have
\[
\cos QX=\frac{x}{\psi AA},\quad\cos QY=\frac{y}{\psi BB},\quad\cos QZ=\frac{z}{\psi CC}
\]
\[
\cos QM=\frac{1}{\psi r}(\frac{(a-x)x}{AA}+\frac{(b-y)y}{BB}+\frac{(c-z)z}{CC})
\]

\Article{12}
\par
Let us now introduce two indeterminates $p,q$ such that
\[
\begin{aligned}
x&=A\cos p\\
y&=B\sin p.\cos q\\
z&=C\sin p.\sin q
\end{aligned}
\]
